\section{Clinical context}
Abdominal aortic aneurysm (AAA) is a localized dilatation of the abdominal aorta, most commonly in its infrarenal segment. The aorta is considered aneurysmal when its diameter exceeds 3~cm or is at least 1.5 times the normal diameter \cite{sakalihasan_abdominal_2018, katsoudas_hemodynamic_2026}. Left untreated, an AAA can rupture, an event associated with high mortality. Elective repair is recommended when the diameter reaches 5.5~cm in men or 5.0~cm in women, when the aneurysm grows rapidly, or when it becomes symptomatic, while smaller aneurysms are monitored with periodic imaging \cite{wanhainen_editors_2024}.
Repair can be performed with open surgery or with endovascular aneurysm repair (EVAR). EVAR is often preferred because it is minimally invasive and is associated with lower perioperative mortality, although this early survival advantage is not maintained in the long term \cite{braak_scoping_2026, wanhainen_editors_2024} \citeneeded{long-term RCTs, e.g. EVAR-1}. Post-procedural complications, however, are common, with an estimated incidence of 18--30\%, and frequently require re-intervention \cite{daye_complications_2018}. % TODO: specify follow-up period of the 18--30% estimate
For this reason, patients treated with EVAR undergo lifelong imaging surveillance, which is costly and exposes them to repeated radiation and iodinated contrast agents \cite{wanhainen_editors_2024}. Accurate prediction of adverse events (AEs) after EVAR could help clinicians decide whether surveillance should be intensified or reduced for individual patients \cite{peng_endoleak_2026} \citeneeded{Miura}.

\section{Post-EVAR adverse events}
% TODO: to be defined with clinicians.
% Planned structure: (1) main complications: endoleaks (types I--V), stent-graft migration,
% limb thrombosis/occlusion, sac enlargement, rupture; (2) clinical criterion used to group
% them into a composite outcome (e.g. events requiring re-intervention); (3) why a composite
% outcome is preferred over a single one such as endoleak.

\section{Risk prediction after EVAR: state of the art and limitations}
Assessing the risk of complications after EVAR remains challenging. Risk scores developed for open repair do not account for the anatomical factors that determine the outcome of endovascular repair \citeneeded{}. Machine learning (ML) methods have therefore been widely explored for post-EVAR risk prediction, as they can handle high-dimensional clinical and imaging data and capture non-linear relationships \cite{braak_scoping_2026, collins_tripodai_2024}.
However, ML-based tools developed specifically for post-EVAR complications still show limited predictive performance. Two limitations recur. First, most models are developed and tested on small, single-center datasets without consistent external validation, which limits their generalizability and clinical applicability \cite{REF19}. Second, they rely mainly on static descriptors extracted from CTA, namely aortic morphology and radiomic features \cite{braak_scoping_2026, REF13, REF20, REF18}, which characterize the shape and tissue composition of the aneurysm but not the flow conditions and mechanical loads acting on the vessel wall and the endograft.
Computational fluid dynamics (CFD) can provide this information through a quantitative description of blood flow patterns \cite{katsoudas_hemodynamic_2026, REF21}. Hemodynamic indicators have shown promise in predicting AAA growth \cite{REF15, REF28, REF12, REF9}, but their relationship with outcomes after EVAR has received little attention.

\section{Research question and objectives}
This thesis addresses the following research question:
\begin{quote}
\emph{Do CFD-derived hemodynamic indicators improve the prediction of post-EVAR complications compared with models based on morphological and radiomic features alone?}
\end{quote}
To answer it, the thesis pursues the following objectives:
\begin{enumerate}
    \item develop a pipeline to extract morphological, radiomic and hemodynamic features from pre-operative CTA scans; % TODO: confirm pre- vs post-operative CTA
    \item train and compare predictive models built on each feature domain and on their combination;
    \item quantify the contribution of hemodynamic indicators to predictive performance and identify the most informative features;
    \item assess the generalizability of the resulting model on an external cohort. % TODO: keep only if external validation is planned
\end{enumerate}